\chapter{Stochastic Processes}\label{martingales}
\chapterstatus{draft}

A \emph{stochastic process} is a family of random variables $(X_t)_{t \in T}$ indexed by time, $T = \ZZ_{\geq 0}$ or $T = [0, \infty)$. Chapter~\ref{probability} studied sums of independent variables; here the
variables may depend on each other, and the chapter treats the three kinds of dependence that dominate probability theory and its applications:
\begin{itemize}
\item \emph{martingales}, the fair games, whose increments have conditional mean zero (Sections~\ref{martingales:section-conditional-expectation}--\ref{martingales:section-convergence});
\item \emph{Markov chains}, processes without memory (Sections~\ref{martingales:section-markov-chains} and~\ref{martingales:section-poisson-process});
\item \emph{Brownian motion}, the continuous-time limit of random walks, and the stochastic calculus built on it (Sections~\ref{martingales:section-brownian-motion} and~\ref{martingales:section-stochastic-calculus}).
\end{itemize}
The basic tool for all three is conditional expectation, which makes ``the expected value of $X$ given the information available at time $n$'' a random variable.

\noindent\textbf{Prerequisites.} Chapter~\ref{probability} (Probability Theory), Chapter~\ref{signed-measures} (Signed and Complex Measures) and Chapter~\ref{functional-analysis} (Functional Analysis).

\section{Conditional expectation}\label{martingales:section-conditional-expectation}

If $B$ is an event with $P(B) > 0$, the conditional probability $P(A \mid B) = P(A \cap B)/P(B)$ and the conditional expectation $E[X \mid B] = E[X; B]/P(B)$ are elementary. If $\Omega$ is partitioned into countably many
events $B_i$ of positive probability, generating a $\sigma$-algebra $\mathcal{G}$, it is natural to call the random variable
\[
E[X \mid \mathcal{G}] = \sum_iE[X \mid B_i]\mathbf{1}_{B_i}
\]
the conditional expectation of $X$ given the information in $\mathcal{G}$: if we know which $B_i$ occurred, it is our best guess of $X$. It is $\mathcal{G}$-measurable, and $E[E[X \mid \mathcal{G}]; A] = E[X; A]$ for every
$A \in \mathcal{G}$. These two properties make sense for every sub-$\sigma$-algebra, and they are taken as the definition.

\begin{definition}\label{martingales:definition-conditional-expectation}
Let $(\Omega, \mathcal{F}, P)$ be a probability space, $\mathcal{G} \subseteq \mathcal{F}$ a sub-$\sigma$-algebra and $X \in L^1(\Omega, \mathcal{F}, P)$. A \emph{conditional expectation} of $X$ given $\mathcal{G}$ is a random variable $Y$
such that
\begin{enumerate}
\item $Y$ is $\mathcal{G}$-measurable and integrable, and
\item $E[Y; A] = E[X; A]$ for every $A \in \mathcal{G}$.
\end{enumerate}
\end{definition}

\begin{notation}\label{martingales:notation-conditional-expectation}
We write $E[X \mid \mathcal{G}]$ for any conditional expectation (they agree almost surely, by Theorem~\ref{martingales:theorem-conditional-expectation}), and all equalities and inequalities between conditional expectations
hold almost surely. $E[X \mid Y] = E[X \mid \sigma(Y)]$ for a random variable $Y$, and $P(A \mid \mathcal{G}) = E[\mathbf{1}_A \mid \mathcal{G}]$.
\end{notation}

\begin{theorem}\label{martingales:theorem-conditional-expectation}
For $X \in L^1$ and a sub-$\sigma$-algebra $\mathcal{G}$, a conditional expectation $E[X \mid \mathcal{G}]$ exists and is unique almost surely. If $X \in L^2$, it is the orthogonal projection of $X$ onto the closed subspace
$L^2(\Omega, \mathcal{G}, P)$ of $L^2(\Omega, \mathcal{F}, P)$.
\end{theorem}

\begin{proof}
\emph{Uniqueness.} If $Y$, $Y^{\prime}$ are both conditional expectations, $A = \{Y > Y^{\prime}\} \in \mathcal{G}$ and $E[Y - Y^{\prime}; A] = E[X; A] - E[X; A] = 0$, so $P(A) = 0$; by symmetry $Y = Y^{\prime}$ almost surely.
\emph{Existence.} Write $X = X^+ - X^-$; it suffices to treat $X \geq 0$. Then $\nu(A) = E[X; A]$, $A \in \mathcal{G}$, is a finite measure on $(\Omega, \mathcal{G})$, absolutely continuous with respect to the restriction of $P$
to $\mathcal{G}$. By the Radon--Nikodym theorem (Theorem~\ref{signed-measures:theorem-lebesgue-radon-nikodym}), $\nu = Y\cdot P|_{\mathcal{G}}$ for a $\mathcal{G}$-measurable $Y \geq 0$, which is a conditional expectation.
\emph{The $L^2$ case.} $L^2(\mathcal{G})$ is complete, hence closed in $L^2(\mathcal{F})$. Let $Y$ be the orthogonal projection of $X$ (Theorem~\ref{functional-analysis:theorem-projection}). Then $X - Y \perp \mathbf{1}_A$ for
$A \in \mathcal{G}$, which is condition (2).
\end{proof}

The $L^2$ description says that $E[X \mid \mathcal{G}]$ is the best $\mathcal{G}$-measurable predictor of $X$ in the mean square sense; this is how conditional expectations appear in filtering and regression.

\begin{proposition}\label{martingales:proposition-conditional-expectation-properties}
Let $X, Y \in L^1$ and $\mathcal{H} \subseteq \mathcal{G} \subseteq \mathcal{F}$.
\begin{enumerate}
\item (Linearity and monotonicity) $E[aX + bY \mid \mathcal{G}] = aE[X \mid \mathcal{G}] + bE[Y \mid \mathcal{G}]$, and $X \leq Y$ implies $E[X \mid \mathcal{G}] \leq E[Y \mid \mathcal{G}]$; so $|E[X \mid \mathcal{G}]| \leq E[|X| \mid \mathcal{G}]$.
\item $E[E[X \mid \mathcal{G}]] = E[X]$.
\item (Taking out what is known) If $Z$ is $\mathcal{G}$-measurable and $ZX$ is integrable, $E[ZX \mid \mathcal{G}] = ZE[X \mid \mathcal{G}]$; in particular $E[Z \mid \mathcal{G}] = Z$.
\item (Tower property) $E[E[X \mid \mathcal{G}] \mid \mathcal{H}] = E[X \mid \mathcal{H}] = E[E[X \mid \mathcal{H}] \mid \mathcal{G}]$.
\item (Independence) If $\sigma(X)$ and $\mathcal{G}$ are independent, $E[X \mid \mathcal{G}] = E[X]$.
\item (Conditional monotone and dominated convergence) If $0 \leq X_n \uparrow X$, then $E[X_n \mid \mathcal{G}] \uparrow E[X \mid \mathcal{G}]$. If $X_n \to X$ almost surely and $|X_n| \leq W \in L^1$, then
$E[X_n \mid \mathcal{G}] \to E[X \mid \mathcal{G}]$ almost surely.
\item (Conditional Jensen) If $\varphi \colon \RR \to \RR$ is convex and $\varphi(X) \in L^1$, then $\varphi(E[X \mid \mathcal{G}]) \leq E[\varphi(X) \mid \mathcal{G}]$. In particular $\|E[X \mid \mathcal{G}]\|_p \leq \|X\|_p$ for $1 \leq p \leq \infty$.
\end{enumerate}
\end{proposition}

\begin{proof}
In each case we check that the proposed right side satisfies Definition~\ref{martingales:definition-conditional-expectation} and use uniqueness.
(1) Linearity is clear. If $X \leq Y$ and $A = \{E[X \mid \mathcal{G}] > E[Y \mid \mathcal{G}]\}$, then $0 \leq E[E[X \mid \mathcal{G}] - E[Y \mid \mathcal{G}]; A] = E[X - Y; A] \leq 0$, so $P(A) = 0$.
(2) Take $A = \Omega$.
(3) For $Z = \mathbf{1}_B$ with $B \in \mathcal{G}$: $E[\mathbf{1}_BE[X \mid \mathcal{G}]; A] = E[X; A \cap B]$. By linearity it holds for simple $Z$; for $X, Z \geq 0$ by monotone convergence (approximate $Z$ by simple functions from below);
in general by splitting into positive and negative parts.
(4) The second equality is (3), since $E[X \mid \mathcal{H}]$ is $\mathcal{G}$-measurable. For the first, if $A \in \mathcal{H} \subseteq \mathcal{G}$, $E[E[X \mid \mathcal{G}]; A] = E[X; A]$.
(5) For $A \in \mathcal{G}$, $E[X; A] = E[X]P(A) = E[E[X]; A]$ by Corollary~\ref{probability:corollary-independence-products}(2).
(6) By (1), $Y_n = E[X_n \mid \mathcal{G}]$ is a.s.\ nondecreasing; its limit $Y$ is $\mathcal{G}$-measurable and $E[Y; A] = \lim E[Y_n; A] = \lim E[X_n; A] = E[X; A]$ by monotone convergence. For the dominated version, let
$V_n = \sup_{k \geq n}|X_k - X|$, so $V_n \downarrow 0$ and $V_n \leq 2W$; then $|E[X_n \mid \mathcal{G}] - E[X \mid \mathcal{G}]| \leq E[V_n \mid \mathcal{G}]$, which is a.s.\ nonincreasing with a limit $V \geq 0$, and $E[V] \leq E[V_n] \to 0$ by
dominated convergence; so $V = 0$.
(7) A convex function is the supremum of the countably many affine functions $\ell(x) = ax + b \leq \varphi$ with rational $a$, $b$ (its supporting lines at rational points can be approximated by these). For each such $\ell$,
$E[\varphi(X) \mid \mathcal{G}] \geq E[\ell(X) \mid \mathcal{G}] = \ell(E[X \mid \mathcal{G}])$ almost surely, and a countable supremum of a.s.\ inequalities holds a.s. With $\varphi(x) = |x|^p$ and (2) we get the $L^p$ bound for $p < \infty$;
for $p = \infty$ use (1).
\end{proof}

\begin{example}\label{martingales:example-conditional-density}
\begin{enumerate}
\item If $(X, Y)$ has a joint density $f(x, y)$ and $g(X)$ is integrable, then
\[
E[g(X) \mid Y] = h(Y), \qquad h(y) = \frac{\int g(x)f(x, y)\,dx}{\int f(x, y)\,dx}
\]
(and $h(y) = 0$ where the denominator vanishes). Indeed, for $A = \{Y \in B\}$, Fubini gives $E[h(Y); A] = \int_B h(y)\int f(x, y)\,dx\,dy = \int_B\int g(x)f(x, y)\,dx\,dy = E[g(X); A]$. The function
$x \mapsto f(x, y)/\int f(x^{\prime}, y)\,dx^{\prime}$ is the \emph{conditional density} of $X$ given $Y = y$.
\item If $X_1, X_2, \ldots$ are i.i.d.\ integrable and $S_n = X_1 + \cdots + X_n$, then $E[X_1 \mid S_n] = S_n/n$: by symmetry $E[X_k \mid S_n]$ is the same for every $k \leq n$ (the pairs $(X_k, S_n)$ all have the same law), and the sum
over $k$ is $E[S_n \mid S_n] = S_n$.
\end{enumerate}
\end{example}

\begin{counterexample}[The Borel--Kolmogorov paradox]\label{martingales:counterexample-borel-kolmogorov}
Conditioning on an event of probability zero is not well defined; only conditioning on a $\sigma$-algebra is. Let a point be uniform on the unit sphere $S^2$, with longitude $\lambda \in [0, 2\pi)$ and latitude
$\theta \in [-\frac\pi2, \frac\pi2]$; its density in these coordinates is $\frac{1}{4\pi}\cos\theta$. Conditioning on the equator $\theta = 0$, by Example~\ref{martingales:example-conditional-density}, the longitude is uniform. But the
equator is also the meridian circle $\lambda \in \{0, \pi\}$ after a rotation, and conditioning on $\lambda$ the latitude has conditional density $\frac12\cos\theta$, not uniform. The same great circle gets two different
``conditional laws'', because it is obtained as the limit of different families of events of positive probability.
\end{counterexample}

\begin{definition}\label{martingales:definition-regular-conditional}
Let $X$ be a real random variable and $\mathcal{G}$ a sub-$\sigma$-algebra. A \emph{regular conditional distribution} of $X$ given $\mathcal{G}$ is a map $\mu \colon \Omega \times \mathcal{B}(\RR) \to [0, 1]$ such that
$\mu(\omega, \cdot)$ is a probability measure for every $\omega$, and $\mu(\cdot, B)$ is a version of $P(X \in B \mid \mathcal{G})$ for every Borel $B$.
\end{definition}

\begin{theorem}\label{martingales:theorem-regular-conditional}
Every real random variable has a regular conditional distribution given any sub-$\sigma$-algebra; then $E[g(X) \mid \mathcal{G}](\omega) = \int g\,d\mu(\omega, \cdot)$ for Borel $g$ with $g(X) \in L^1$.
\end{theorem}

\begin{proof}[Sketch of proof]
For rational $q$ choose versions $F(\omega, q)$ of $P(X \leq q \mid \mathcal{G})$. By Proposition~\ref{martingales:proposition-conditional-expectation-properties}(1) and (6), outside a null set $N$ the function $q \mapsto F(\omega, q)$
is nondecreasing, $F(\omega, q) = \lim_{r \downarrow q}F(\omega, r)$ along rationals, and its limits at $\mp\infty$ are $0$ and $1$ (countably many conditions). For $\omega \notin N$ let $F(\omega, x) = \inf_{q > x}F(\omega, q)$, a distribution
function, and $\mu(\omega, \cdot)$ the corresponding law (Proposition~\ref{probability:proposition-distribution-function}); for $\omega \in N$ let $\mu(\omega, \cdot)$ be a fixed law. The class of Borel $B$ for which $\mu(\cdot, B)$ is a
version of $P(X \in B \mid \mathcal{G})$ contains the half-lines and is a Dynkin system, by conditional monotone convergence, hence is everything (Lemma~\ref{measure-theory:lemma-dynkin}). The formula for $g$ follows for
indicators, then simple functions, then by monotone convergence. The verifications omitted are the monotone limits at each step.
\end{proof}

The same holds for random variables with values in any complete separable metric space, which is the setting in which regular conditional distributions are used \cite{kallenberg}; for arbitrary measurable spaces it
can fail.

\section{Martingales and stopping times}\label{martingales:section-martingales}

\begin{definition}\label{martingales:definition-martingale}
A \emph{filtration} is a nondecreasing sequence $\mathcal{F}_0 \subseteq \mathcal{F}_1 \subseteq \cdots$ of sub-$\sigma$-algebras of $\mathcal{F}$; $\mathcal{F}_n$ is the information available at time $n$. A process $(X_n)_{n \geq 0}$ is
\emph{adapted} if each $X_n$ is $\mathcal{F}_n$-measurable, and \emph{predictable} if each $X_n$ ($n \geq 1$) is $\mathcal{F}_{n-1}$-measurable. An adapted process of integrable random variables is a
\begin{enumerate}
\item \emph{martingale} if $E[X_{n+1} \mid \mathcal{F}_n] = X_n$ for all $n$;
\item \emph{submartingale} if $E[X_{n+1} \mid \mathcal{F}_n] \geq X_n$, and \emph{supermartingale} if $E[X_{n+1} \mid \mathcal{F}_n] \leq X_n$.
\end{enumerate}
Without mention of a filtration, the \emph{natural filtration} $\mathcal{F}_n = \sigma(X_0, \ldots, X_n)$ is meant.
\end{definition}

\begin{notation}\label{martingales:notation-martingale}
$\Delta X_n = X_n - X_{n-1}$ denotes the increments, and $\mathcal{F}_\infty = \sigma(\bigcup_n\mathcal{F}_n)$. By the tower property, a martingale satisfies $E[X_m \mid \mathcal{F}_n] = X_n$ for all $m \geq n$, and $E[X_n] = E[X_0]$.
\end{notation}

A martingale is the fortune of a gambler in a fair game: given the past, the expected gain of the next round is zero. A supermartingale is an unfavourable game (the name comes from potential theory: superharmonic functions
are supermartingales along Brownian paths).

\begin{example}\label{martingales:example-martingales}
\begin{enumerate}
\item \emph{Sums.} If $\xi_1, \xi_2, \ldots$ are independent integrable with mean $0$, $S_n = \xi_1 + \cdots + \xi_n$ is a martingale: $E[S_{n+1} \mid \mathcal{F}_n] = S_n + E[\xi_{n+1}] = S_n$ by
Proposition~\ref{martingales:proposition-conditional-expectation-properties}(3), (5). If moreover $E[\xi_k^2] = \sigma^2$, then $S_n^2 - n\sigma^2$ is a martingale.
\item \emph{Products.} If $\xi_k \geq 0$ are independent with mean $1$, $M_n = \xi_1\cdots\xi_n$ is a martingale. For instance, for i.i.d.\ $\xi_k$ with $\Lambda(\lambda) = \log E[e^{\lambda\xi_1}] < \infty$,
$\exp(\lambda S_n - n\Lambda(\lambda))$ is a martingale (the \emph{exponential martingale}).
\item \emph{Doob martingales.} If $X \in L^1$, $X_n = E[X \mid \mathcal{F}_n]$ is a martingale, by the tower property: the successive best guesses of $X$ as information accumulates.
\item \emph{Likelihood ratios.} If $P$ and $Q$ are probability measures with $Q \ll P$ on each $\mathcal{F}_n$, the densities $L_n = dQ|_{\mathcal{F}_n}/dP|_{\mathcal{F}_n}$ form a $P$-martingale.
\item \emph{Convex functions.} If $X_n$ is a martingale and $\varphi$ is convex with $\varphi(X_n) \in L^1$, then $\varphi(X_n)$ is a submartingale, by conditional Jensen; for instance $|X_n|$, and $X_n^2$ if $X_n \in L^2$.
\end{enumerate}
\end{example}

\begin{definition}\label{martingales:definition-transform}
For an adapted process $X$ and a predictable process $H$, the \emph{martingale transform} is $(H \cdot X)_n = \sum_{k=1}^nH_k\,\Delta X_k$, with $(H \cdot X)_0 = 0$.
\end{definition}

If $X_n$ is the price of a share and $H_k$ the number of shares held during the round $k$, chosen with the information at time $k - 1$, then $(H \cdot X)_n$ is the gain of the trading strategy $H$. The discrete analogue of the
stochastic integral of Section~\ref{martingales:section-stochastic-calculus}, it is a martingale whenever $X$ is: no strategy beats a fair game.

\begin{proposition}\label{martingales:proposition-transform}
If $X$ is a martingale and $H$ is predictable and bounded, $H \cdot X$ is a martingale. If $X$ is a supermartingale and $0 \leq H_k \leq c$, $H \cdot X$ is a supermartingale.
\end{proposition}

\begin{proof}
$(H \cdot X)_n$ is integrable and adapted, and $E[(H \cdot X)_{n+1} - (H \cdot X)_n \mid \mathcal{F}_n] = H_{n+1}E[X_{n+1} - X_n \mid \mathcal{F}_n]$, which is $0$, resp.\ $\leq 0$, by
Proposition~\ref{martingales:proposition-conditional-expectation-properties}(3).
\end{proof}

\begin{definition}\label{martingales:definition-stopping-time}
A \emph{stopping time} is a random variable $T \colon \Omega \to \{0, 1, 2, \ldots, \infty\}$ with $\{T = n\} \in \mathcal{F}_n$ for every $n$: whether to stop at time $n$ is decided with the information at time $n$. The
\emph{stopped process} is $X^T_n = X_{T \wedge n}$; if $T < \infty$, $X_T(\omega) = X_{T(\omega)}(\omega)$.
\end{definition}

\begin{example}\label{martingales:example-stopping-times}
Constants are stopping times, and so are $T \wedge S$, $T \vee S$ and $T + S$ for stopping times $T$, $S$. For an adapted process and a Borel set $B$, the \emph{hitting time} $T_B = \inf\{n \geq 0 : X_n \in B\}$ is a stopping time,
since $\{T_B = n\} = \{X_0 \notin B, \ldots, X_{n-1} \notin B, X_n \in B\}$. The \emph{last} time $X$ is in $B$ before time $10$ is not a stopping time: it depends on the future.
\end{example}

\begin{proposition}\label{martingales:proposition-stopped}
If $X$ is a martingale (resp.\ supermartingale) and $T$ a stopping time, $X^T$ is a martingale (resp.\ supermartingale). If $T \leq N$ is bounded, $E[X_T] = E[X_0]$ (resp.\ $E[X_T] \leq E[X_0]$).
\end{proposition}

\begin{proof}
$X_{T \wedge n} = X_0 + (H \cdot X)_n$ with $H_k = \mathbf{1}_{\{T \geq k\}}$, which is predictable since $\{T \geq k\} = \Omega \setminus \{T \leq k - 1\} \in \mathcal{F}_{k-1}$. Apply Proposition~\ref{martingales:proposition-transform}, and
for bounded $T$ take $n = N$.
\end{proof}

\begin{theorem}[Optional stopping theorem]\label{martingales:theorem-optional-stopping}
Let $X$ be a martingale and $T$ an almost surely finite stopping time. Then $E[X_T] = E[X_0]$ in each of the following cases:
\begin{enumerate}
\item $T$ is bounded;
\item $|X_{T \wedge n}| \leq K$ for all $n$, for a constant $K$;
\item $E[T] < \infty$ and $|\Delta X_n| \leq c$ for all $n$, for a constant $c$.
\end{enumerate}
More generally, it holds whenever the family $(X_{T \wedge n})_n$ is uniformly integrable.
\end{theorem}

\begin{proof}
$X_{T \wedge n} \to X_T$ almost surely as $T < \infty$, and $E[X_{T \wedge n}] = E[X_0]$ by Proposition~\ref{martingales:proposition-stopped}. In case (2), dominated convergence gives the claim; in case (3), $|X_{T \wedge n}| \leq |X_0| + cT$, which
is integrable. In the uniformly integrable case use Vitali's theorem (Theorem~\ref{signed-measures:theorem-vitali-convergence}).
\end{proof}

\begin{example}[Gambler's ruin]\label{martingales:example-gamblers-ruin}
A gambler starts with $a$ euros and bets one euro on fair coin tosses until reaching $N > a$ or $0$. The fortune is $S_n = a + \xi_1 + \cdots + \xi_n$ with fair signs $\xi_k$, stopped at $T = \inf\{n : S_n \in \{0, N\}\}$.
\begin{enumerate}
\item $T < \infty$ almost surely, and even $E[T] < \infty$: in every block of $N$ consecutive tosses, all heads (probability $2^{-N}$) forces exit, so $P(T > kN) \leq (1 - 2^{-N})^k$.
\item By optional stopping (case (2)), $a = E[S_T] = NP(S_T = N)$: the probability of winning is $a/N$.
\item $S_n^2 - n$ is a martingale with bounded increments on $\{n < T\}$; optional stopping applied to it (case (3) for the stopped process) gives $E[S_T^2] = a^2 + E[T]$, so $E[T] = N^2\frac aN - a^2 = a(N - a)$.
\item For an unfair coin, $P(\xi_k = 1) = p \neq \frac12$, $q = 1 - p$, the process $(q/p)^{S_n}$ is a martingale, and the same argument gives $P(S_T = N) = \frac{1 - (q/p)^a}{1 - (q/p)^N}$. In a casino with $p = \frac{18}{37}$, starting with
$a = 100$ and aiming at $N = 200$, the probability of success is about $4.4 \cdot 10^{-3}$, while in a fair game it is $\frac12$.
\end{enumerate}
\end{example}

\begin{counterexample}\label{martingales:counterexample-optional-stopping}
Optional stopping fails without conditions. For the simple random walk $S_n$ started at $0$ and $T = \inf\{n : S_n = 1\}$, Example~\ref{martingales:example-gamblers-ruin} shows that $P(T < \infty) \geq P(\text{hit } 1 \text{ before } {-N})
= \frac{N}{N + 1}$ for every $N$, so $T < \infty$ almost surely; but $E[S_T] = 1 \neq 0 = E[S_0]$. By case (3) of the theorem, $E[T] = \infty$. This is the \emph{doubling strategy} in disguise: a gambler who doubles the stake after
each loss is sure to win one euro eventually, but needs unbounded credit and time.
\end{counterexample}

\begin{theorem}[Wald's identity]\label{martingales:theorem-wald}
Let $\xi_1, \xi_2, \ldots$ be i.i.d.\ integrable with mean $m$, $S_n = \sum_{k \leq n}\xi_k$, and $T$ a stopping time for the natural filtration with $E[T] < \infty$. Then $E[S_T] = mE[T]$.
\end{theorem}

\begin{proof}
$M_n = S_n - nm$ is a martingale, so $E[M_{T \wedge n}] = 0$. Now $|M_{T \wedge n}| \leq \sum_{k=1}^\infty|\xi_k - m|\mathbf{1}_{\{T \geq k\}}$, and since $\{T \geq k\} \in \mathcal{F}_{k-1}$ is independent of $\xi_k$,
$E\sum_k|\xi_k - m|\mathbf{1}_{\{T \geq k\}} = E|\xi_1 - m|\sum_kP(T \geq k) = E|\xi_1 - m|E[T] < \infty$. Dominated convergence gives $E[S_T - Tm] = 0$.
\end{proof}

\section{Martingale inequalities and convergence}\label{martingales:section-convergence}

\begin{theorem}[Doob's maximal inequality]\label{martingales:theorem-doob-maximal}
Let $X$ be a submartingale, $\lambda > 0$ and $X_n^* = \max_{k \leq n}X_k$. Then $\lambda P(X_n^* \geq \lambda) \leq E[X_n; X_n^* \geq \lambda] \leq E[X_n^+]$.
\end{theorem}

\begin{proof}
Let $A_k = \{X_k \geq \lambda, X_j < \lambda \text{ for } j < k\} \in \mathcal{F}_k$, disjoint with union $\{X_n^* \geq \lambda\}$. By the submartingale property $E[X_n; A_k] \geq E[X_k; A_k] \geq \lambda P(A_k)$; sum over $k \leq n$.
\end{proof}

Applied to the submartingale $S_n^2$ of a sum of independent centred variables, this is Kolmogorov's inequality (Theorem~\ref{probability:theorem-kolmogorov-inequality}).

\begin{theorem}[Doob's $L^p$ inequality]\label{martingales:theorem-doob-lp}
If $X$ is a nonnegative submartingale and $1 < p < \infty$, then $\|X_n^*\|_p \leq \frac{p}{p - 1}\|X_n\|_p$. In particular, for a martingale $X$, $E[\max_{k \leq n}|X_k|^2] \leq 4E[X_n^2]$.
\end{theorem}

\begin{proof}
We may assume $X_n \in L^p$; then $X_k \in L^p$ for $k \leq n$ by conditional Jensen ($X_k \leq E[X_n \mid \mathcal{F}_k]$), so $X_n^* \leq \sum_{k \leq n}X_k \in L^p$. By the layer cake formula
(Proposition~\ref{measure-theory:proposition-layer-cake}), Theorem~\ref{martingales:theorem-doob-maximal} and Fubini,
\[
\begin{aligned}
E[(X_n^*)^p] &= \int_0^\infty p\lambda^{p-1}P(X_n^* \geq \lambda)\,d\lambda \leq \int_0^\infty p\lambda^{p-2}E[X_n; X_n^* \geq \lambda]\,d\lambda \\
&= \frac{p}{p - 1}E[X_n(X_n^*)^{p-1}].
\end{aligned}
\]
By H\"older's inequality the right side is at most $\frac{p}{p - 1}\|X_n\|_p\|X_n^*\|_p^{p-1}$; divide by $\|X_n^*\|_p^{p-1}$ (if it is nonzero). For a martingale apply this to $|X_n|$.
\end{proof}

\begin{counterexample}\label{martingales:counterexample-doob-l1}
There is no $L^1$ version: for the martingale $X_n = \prod_{k \leq n}\xi_k$ with independent $\xi_k \in \{0, 2\}$ equally likely, $E[X_n] = 1$, while $X_n^* = 2^{N}$ where $N = \min(n, \text{the index of the first } 0)$, so
$E[X_n^*] = \sum_{k<n}2^k2^{-k-1} + 2^n2^{-n} = \frac n2 + 1 \to \infty$.
\end{counterexample}

\begin{definition}\label{martingales:definition-upcrossings}
For a process $X$, numbers $a < b$ and $n \geq 0$, the \emph{number of upcrossings} $U_n[a, b]$ of $[a, b]$ by time $n$ is the largest $k$ for which there are times $s_1 < t_1 < s_2 < t_2 < \cdots < s_k < t_k \leq n$ with
$X_{s_i} \leq a$ and $X_{t_i} \geq b$.
\end{definition}

\begin{lemma}[Doob's upcrossing inequality]\label{martingales:lemma-upcrossing}
If $X$ is a submartingale, $(b - a)E[U_n[a, b]] \leq E[(X_n - a)^+] - E[(X_0 - a)^+]$.
\end{lemma}

\begin{proof}
$Y_n = a + (X_n - a)^+$ is a submartingale (Example~\ref{martingales:example-martingales}(5), as $x \mapsto a + (x - a)^+$ is convex and nondecreasing) with the same upcrossings of $[a, b]$. Let $H_k = 1$ if time $k$ lies in an
upcrossing in progress, precisely: $H$ switches on after the first time $Y \leq a$ (i.e.\ $Y = a$) and off after the next time $Y \geq b$, and so on; $H$ is predictable, since $H_k$ is decided by $Y_0, \ldots, Y_{k-1}$. Each completed
upcrossing contributes at least $b - a$ to $(H \cdot Y)_n$, and an incomplete final one contributes $Y_n - a \geq 0$; so $(b - a)U_n[a, b] \leq (H \cdot Y)_n$. As $1 - H$ is predictable and nonnegative, $((1 - H) \cdot Y)_n$ is a
submartingale with mean $\geq 0$ (Proposition~\ref{martingales:proposition-transform} applied to the supermartingale $-Y$). Hence $(b - a)E[U_n] \leq E[(H \cdot Y)_n] \leq E[(H \cdot Y)_n] + E[((1 - H) \cdot Y)_n] = E[Y_n - Y_0]$.
\end{proof}

\begin{theorem}[Martingale convergence theorem]\label{martingales:theorem-convergence}
Let $X$ be a submartingale with $\sup_nE[X_n^+] < \infty$. Then $X_n$ converges almost surely to an integrable random variable $X_\infty$. In particular, every martingale bounded in $L^1$ and every nonnegative
supermartingale converges almost surely.
\end{theorem}

\begin{proof}
For rationals $a < b$, $U_n[a, b] \uparrow U_\infty[a, b]$ and by Lemma~\ref{martingales:lemma-upcrossing} and monotone convergence $E[U_\infty[a, b]] \leq \frac{\sup_nE[X_n^+] + |a|}{b - a} < \infty$, so $U_\infty[a, b] < \infty$ almost
surely. Outside the countable union of the null sets $\{U_\infty[a, b] = \infty\}$, $\liminf X_n = \limsup X_n$: if $\liminf X_n < a < b < \limsup X_n$ there would be infinitely many upcrossings. So $X_n \to X_\infty \in [-\infty, \infty]$
almost surely, and by Fatou's lemma (Corollary~\ref{measure-theory:corollary-integral-properties}) $E|X_\infty| \leq \liminf E|X_n| = \liminf(2E[X_n^+] - E[X_n]) \leq 2\sup E[X_n^+] - E[X_0] < \infty$. For a nonnegative
supermartingale apply this to the submartingale $-X_n$, whose positive part vanishes.
\end{proof}

\begin{example}\label{martingales:example-convergence}
\begin{enumerate}
\item \emph{P\'olya's urn.} An urn contains one red and one blue ball; at each step a ball is drawn uniformly, and returned together with a new ball of the same colour. The fraction $R_n$ of red balls after $n$ steps (out
of $n + 2$) is a martingale with values in $[0, 1]$: if there are $r$ red balls, $E[R_{n+1} \mid \mathcal{F}_n] = \frac{r}{n+2}\cdot\frac{r + 1}{n + 3} + \frac{n + 2 - r}{n + 2}\cdot\frac{r}{n + 3} = \frac{r}{n + 2}$. So $R_n$ converges almost
surely; the limit is uniformly distributed on $[0, 1]$ (Exercise~\ref{martingales:exercise-polya}). The urn ``remembers'' its early history forever.
\item \emph{Random signs.} $\sum_k\varepsilon_k/k$ with independent fair signs is an $L^2$-bounded martingale, so it converges almost surely, recovering Example~\ref{probability:example-random-harmonic}.
\end{enumerate}
\end{example}

\begin{counterexample}\label{martingales:counterexample-not-l1}
Almost sure convergence need not hold in $L^1$: the martingale $X_n = \prod_{k \leq n}\xi_k$ of Counterexample~\ref{martingales:counterexample-doob-l1} converges almost surely to $0$ (some $\xi_k$ vanishes), while
$E[X_n] = 1$ for all $n$. A fair game can converge to ruin with probability one.
\end{counterexample}

\begin{theorem}\label{martingales:theorem-ui-convergence}
For a martingale $(X_n)$ the following are equivalent:
\begin{enumerate}
\item $(X_n)$ is uniformly integrable (Definition~\ref{signed-measures:definition-uniformly-integrable});
\item $X_n$ converges almost surely and in $L^1$;
\item $X_n$ converges in $L^1$;
\item there is $X \in L^1$ with $X_n = E[X \mid \mathcal{F}_n]$ for all $n$ (the martingale is \emph{closed}).
\end{enumerate}
Then $X_n \to E[X \mid \mathcal{F}_\infty]$, and $X_\infty = \lim X_n$ satisfies $X_n = E[X_\infty \mid \mathcal{F}_n]$.
\end{theorem}

\begin{proof}
(1) $\Rightarrow$ (2): A uniformly integrable family is bounded in $L^1$ (Proposition~\ref{signed-measures:proposition-uniformly-integrable}), so $X_n \to X_\infty$ almost surely by Theorem~\ref{martingales:theorem-convergence}, and in $L^1$ by
Vitali's theorem. (2) $\Rightarrow$ (3) is clear. (3) $\Rightarrow$ (4): If $X_n \to X_\infty$ in $L^1$, then for fixed $n$ and $m \geq n$, $X_n = E[X_m \mid \mathcal{F}_n]$, and $\|E[X_m \mid \mathcal{F}_n] - E[X_\infty \mid \mathcal{F}_n]\|_1 \leq
\|X_m - X_\infty\|_1 \to 0$ (conditional Jensen); so $X_n = E[X_\infty \mid \mathcal{F}_n]$.
(4) $\Rightarrow$ (1): Let $Y_n = E[X \mid \mathcal{F}_n]$. Given $\varepsilon > 0$, choose $\delta$ with $E[|X|; A] < \varepsilon$ when $P(A) < \delta$ (Proposition~\ref{signed-measures:proposition-uniformly-integrable} for the single function $X$). By
Markov's inequality, $P(|Y_n| > M) \leq E|Y_n|/M \leq E|X|/M < \delta$ for $M$ large, uniformly in $n$; and since $\{|Y_n| > M\} \in \mathcal{F}_n$, $E[|Y_n|; |Y_n| > M] \leq E[E[|X| \mid \mathcal{F}_n]; |Y_n| > M] = E[|X|; |Y_n| > M] <
\varepsilon$.
Finally, if $X_n = E[X \mid \mathcal{F}_n] \to X_\infty$, then $X_\infty$ is $\mathcal{F}_\infty$-measurable, and $E[X_\infty; A] = \lim E[X_n; A] = E[X; A]$ for $A \in \bigcup_n\mathcal{F}_n$. The class of $A \in \mathcal{F}_\infty$ with $E[X_\infty; A] = E[X; A]$ contains the
$\pi$-system $\bigcup_n\mathcal{F}_n$, which contains $\Omega$, and is closed under complements and (by dominated convergence) countable disjoint unions; by Lemma~\ref{measure-theory:lemma-dynkin} it is $\mathcal{F}_\infty$, so
$X_\infty = E[X \mid \mathcal{F}_\infty]$.
\end{proof}

\begin{corollary}[L\'evy's upward theorem and zero-one law]\label{martingales:corollary-levy-upward}
For $X \in L^1$, $E[X \mid \mathcal{F}_n] \to E[X \mid \mathcal{F}_\infty]$ almost surely and in $L^1$. In particular, for $A \in \mathcal{F}_\infty$, $P(A \mid \mathcal{F}_n) \to \mathbf{1}_A$ almost surely.
\end{corollary}

\begin{proof}
This is Theorem~\ref{martingales:theorem-ui-convergence}, (4) $\Rightarrow$ (2), with the identification of the limit; for $A \in \mathcal{F}_\infty$, $E[\mathbf{1}_A \mid \mathcal{F}_\infty] = \mathbf{1}_A$.
\end{proof}

L\'evy's zero-one law gives a second proof of Kolmogorov's zero-one law (Theorem~\ref{probability:theorem-zero-one}): if $A$ is a tail event of independent $X_k$ and $\mathcal{F}_n = \sigma(X_1, \ldots, X_n)$, then $A$ is
independent of $\mathcal{F}_n$, so $P(A) = P(A \mid \mathcal{F}_n) \to \mathbf{1}_A$, and $P(A) \in \{0, 1\}$.

\begin{theorem}[Backward martingales]\label{martingales:theorem-backward}
Let $\mathcal{G}_0 \supseteq \mathcal{G}_1 \supseteq \cdots$ be a nonincreasing sequence of sub-$\sigma$-algebras, $\mathcal{G}_\infty = \bigcap_n\mathcal{G}_n$, and $X \in L^1$. Then $E[X \mid \mathcal{G}_n] \to E[X \mid \mathcal{G}_\infty]$
almost surely and in $L^1$.
\end{theorem}

\begin{proof}
Let $Y_n = E[X \mid \mathcal{G}_n]$. For fixed $N$, the sequence $Y_N, Y_{N-1}, \ldots, Y_0$ is a martingale for the filtration $\mathcal{G}_N \subseteq \cdots \subseteq \mathcal{G}_0$, so by Lemma~\ref{martingales:lemma-upcrossing} the expected
number of upcrossings of $[a, b]$ by $(Y_N, \ldots, Y_0)$ is at most $\frac{E|X| + |a|}{b - a}$, uniformly in $N$. As in Theorem~\ref{martingales:theorem-convergence}, $Y_n$ converges almost surely to some $Y_\infty$. The family $(Y_n)$
is uniformly integrable by the argument of (4) $\Rightarrow$ (1) above, so the convergence also holds in $L^1$. $Y_\infty$ is $\mathcal{G}_m$-measurable for every $m$ (as a limit of $Y_n$, $n \geq m$), hence
$\mathcal{G}_\infty$-measurable, and for $A \in \mathcal{G}_\infty$, $E[Y_\infty; A] = \lim E[Y_n; A] = E[X; A]$.
\end{proof}

\begin{corollary}[Strong law of large numbers]\label{martingales:corollary-strong-law}
If $X_1, X_2, \ldots$ are i.i.d.\ and integrable, $S_n/n \to E[X_1]$ almost surely and in $L^1$.
\end{corollary}

\begin{proof}
Let $\mathcal{G}_n = \sigma(S_n, X_{n+1}, X_{n+2}, \ldots)$, a nonincreasing sequence. As in Example~\ref{martingales:example-conditional-density}(2), $E[X_k \mid \mathcal{G}_n]$ is the same for all $k \leq n$: for
$A = \{S_n \in B, (X_{n+1}, X_{n+2}, \ldots) \in C\}$, $E[X_k; A]$ does not depend on $k \leq n$, because the law of $(X_k, S_n, X_{n+1}, \ldots)$ does not, and such sets form a $\pi$-system generating $\mathcal{G}_n$. Hence
$E[X_1 \mid \mathcal{G}_n] = \frac1n\sum_{k \leq n}E[X_k \mid \mathcal{G}_n] = \frac{S_n}{n}$. By Theorem~\ref{martingales:theorem-backward}, $S_n/n$ converges almost surely and in $L^1$. The limit equals $\limsup_n(X_m + \cdots + X_n)/n$ for every $m$,
so it is measurable for the tail $\sigma$-algebra of the $X_k$, hence almost surely constant (Theorem~\ref{probability:theorem-zero-one}); by $L^1$ convergence the constant is $E[S_n/n] = E[X_1]$.
\end{proof}

This is a second proof of Theorem~\ref{probability:theorem-strong-law} for i.i.d.\ variables, adding the $L^1$ convergence. The same idea, conditioning on the symmetric information, proves de Finetti's theorem: an
infinite exchangeable sequence is a mixture of i.i.d.\ sequences \cite{kallenberg}.

\begin{theorem}[Azuma--Hoeffding inequality]\label{martingales:theorem-azuma}
Let $X$ be a martingale with $|\Delta X_k| \leq c_k$. Then for $t > 0$,
\[
P(X_n - X_0 \geq t) \leq \exp\Bigl(-\frac{t^2}{2\sum_{k \leq n}c_k^2}\Bigr).
\]
\end{theorem}

\begin{proof}
By convexity, for $|x| \leq c$, $e^{\lambda x} \leq \frac{c - x}{2c}e^{-\lambda c} + \frac{c + x}{2c}e^{\lambda c}$. Applied to $x = \Delta X_k$ and conditioned on $\mathcal{F}_{k-1}$, where $E[\Delta X_k \mid \mathcal{F}_{k-1}] = 0$, this gives
$E[e^{\lambda\Delta X_k} \mid \mathcal{F}_{k-1}] \leq \cosh(\lambda c_k) \leq e^{\lambda^2c_k^2/2}$ (compare the Taylor series termwise). By induction, using Proposition~\ref{martingales:proposition-conditional-expectation-properties}(3),
$E[e^{\lambda(X_n - X_0)}] \leq e^{\lambda^2\sum c_k^2/2}$. The Chernoff bound (Proposition~\ref{probability:proposition-chernoff}) with $\lambda = t/\sum c_k^2$ concludes.
\end{proof}

\begin{corollary}[Bounded differences]\label{martingales:corollary-bounded-differences}
Let $\xi_1, \ldots, \xi_n$ be independent and $f$ a function of $n$ variables such that changing the $k$-th variable changes $f$ by at most $c_k$. Then $Z = f(\xi_1, \ldots, \xi_n)$ satisfies
$P(Z - EZ \geq t) \leq e^{-t^2/2\sum c_k^2}$.
\end{corollary}

\begin{proof}
Apply Theorem~\ref{martingales:theorem-azuma} to the Doob martingale $X_k = E[Z \mid \xi_1, \ldots, \xi_k]$. By independence, $X_k = \int f(\xi_1, \ldots, \xi_k, y_{k+1}, \ldots, y_n)\,d\mu(y)$ with $\mu$ the law of $(\xi_{k+1}, \ldots, \xi_n)$ (Fubini), and
$X_{k-1}$ is the same integral with $\xi_k$ replaced by an integration variable; so $|\Delta X_k| \leq c_k$.
\end{proof}

\begin{example}\label{martingales:example-bounded-differences}
Throw $m$ balls independently into $n$ bins and let $Z$ be the number of empty bins. Moving one ball changes $Z$ by at most $1$, so $P(|Z - EZ| \geq t) \leq 2e^{-t^2/2m}$, with $E[Z] = n(1 - \frac1n)^m$. For $m = n$, $Z$ is
$n/e + O(\sqrt{n\log n})$ with high probability. The same argument shows that the chromatic number of a random graph is concentrated in an interval of length $O(\sqrt n)$ (Chapter~\ref{probabilistic-method}).
\end{example}

\begin{definition}\label{martingales:definition-branching}
A \emph{Galton--Watson branching process} with offspring law $(p_k)_{k \geq 0}$ is $Z_0 = 1$, $Z_{n+1} = \sum_{i=1}^{Z_n}\xi_{n+1,i}$, where the $\xi_{n,i}$ are i.i.d.\ with $P(\xi = k) = p_k$: each individual of generation $n$
has an independent random number of children. Let $m = E[\xi] = \sum kp_k$ and $f(s) = \sum_kp_ks^k$ the generating function.
\end{definition}

\begin{theorem}\label{martingales:theorem-branching}
Assume $0 < m < \infty$ and $p_1 < 1$. Then $Z_n/m^n$ is a nonnegative martingale, so it converges almost surely. The \emph{extinction probability} $q = P(Z_n = 0 \text{ for some } n)$ is the smallest root of $f(s) = s$ in
$[0, 1]$; in particular $q = 1$ if $m \leq 1$ and $q < 1$ if $m > 1$.
\end{theorem}

\begin{proof}
\emph{Martingale.} Given $\mathcal{F}_n$, $Z_{n+1}$ is a sum of $Z_n$ independent copies of $\xi$, so $E[Z_{n+1} \mid \mathcal{F}_n] = mZ_n$ (decompose on the events $\{Z_n = j\}$). Theorem~\ref{martingales:theorem-convergence} applies.
\emph{Generating functions.} Similarly $E[s^{Z_{n+1}} \mid \mathcal{F}_n] = f(s)^{Z_n}$, so $f_n(s) = E[s^{Z_n}]$ satisfies $f_{n+1} = f_n \circ f$ and $f_n = f^{\circ n}$. Thus $q_n = P(Z_n = 0) = f_n(0)$ satisfies $q_{n+1} = f(q_n)$,
$q_0 = 0$, and $q_n \uparrow q$ (extinction by time $n$ is an increasing union). By continuity $f(q) = q$. If $s \in [0, 1]$ is any fixed point, $q_n = f^{\circ n}(0) \leq f^{\circ n}(s) = s$ by monotonicity of $f$, so $q \leq s$.
\emph{Location of the root.} $f$ is convex on $[0, 1]$ with $f(1) = 1$ and $f^{\prime}(1^-) = m$. If $m \leq 1$ and $p_0 + p_1 < 1$, $f$ is strictly convex, so $f(s) > 1 + m(s - 1) \geq s$ for $s < 1$, and $q = 1$; if $p_0 + p_1 = 1$ and
$m \leq 1$, then $p_0 > 0$ (as $p_1 < 1$) and $f(s) = p_0 + p_1s > s$ for $s < 1$. If $m > 1$, $f(s) < s$ for $s$ slightly below $1$ and $f(0) = p_0 \geq 0$, so there is a fixed point in $[0, 1)$, and $q < 1$.
\end{proof}

\begin{example}\label{martingales:example-branching}
\begin{enumerate}
\item If each individual has $0$ or $2$ children with probabilities $\frac14$, $\frac34$, then $m = \frac32$ and $f(s) = \frac14 + \frac34s^2 = s$ has roots $\frac13$ and $1$: the population survives forever with probability $\frac23$.
\item In the critical case $m = 1$, $Z_n$ is a martingale that converges almost surely to $0$ although $E[Z_n] = 1$: another fair game with certain ruin. Galton and Watson studied the problem in 1874 to explain the
extinction of aristocratic surnames.
\end{enumerate}
\end{example}

\begin{remark}\label{martingales:remark-radon-nikodym}
Martingales also prove theorems outside probability. If $\nu \ll \mu$ are finite measures on $[0, 1)$ and $\mathcal{F}_n$ is generated by the dyadic intervals of length $2^{-n}$, then
$X_n = \sum_I\frac{\nu(I)}{\mu(I)}\mathbf{1}_I$ is a $\mu$-martingale (Example~\ref{martingales:example-martingales}(4)), uniformly integrable because $\nu \ll \mu$; its limit is $d\nu/d\mu$, which proves the Radon--Nikodym theorem
for separable $\sigma$-algebras and the differentiation theorem along dyadic intervals simultaneously \cite{williams-martingales}.
\end{remark}

\section{Markov chains}\label{martingales:section-markov-chains}

\begin{definition}\label{martingales:definition-markov-chain}
Let $S$ be a countable set (the \emph{state space}). A \emph{transition matrix} is a matrix $P = (p(i, j))_{i,j \in S}$ with $p(i, j) \geq 0$ and $\sum_jp(i, j) = 1$. A process $(X_n)_{n \geq 0}$ with values in $S$ is a
\emph{Markov chain} with transition matrix $P$ if
\[
P(X_{n+1} = j \mid X_0 = i_0, \ldots, X_{n-1} = i_{n-1}, X_n = i) = p(i, j)
\]
whenever the conditioning event has positive probability. The law of $X_0$ is the \emph{initial distribution}.
\end{definition}

\begin{notation}\label{martingales:notation-markov}
$P_i$ and $E_i$ denote probability and expectation for the chain started at $X_0 = i$, and $P_\mu$ for the initial distribution $\mu$. Measures on $S$ are row vectors and functions are column vectors: $(\mu P)(j) = \sum_i\mu(i)p(i, j)$,
$(Pf)(i) = \sum_jp(i, j)f(j)$. $p^n(i, j)$ denotes the entries of $P^n$. $T_j = \inf\{n \geq 1 : X_n = j\}$ is the first return or hitting time of $j$.
\end{notation}

\begin{proposition}\label{martingales:proposition-markov-construction}
For every transition matrix $P$ and initial distribution $\mu$ there is a Markov chain with these data. Then $P_\mu(X_0 = i_0, \ldots, X_n = i_n) = \mu(i_0)p(i_0, i_1)\cdots p(i_{n-1}, i_n)$, the law of $X_n$ is $\mu P^n$, and
(\emph{Chapman--Kolmogorov}) $p^{m+n}(i, j) = \sum_kp^m(i, k)p^n(k, j)$.
\end{proposition}

\begin{proof}
Enumerate $S = \{s_1, s_2, \ldots\}$ and let $\phi(i, u) = s_m$ for the least $m$ with $\sum_{l \leq m}p(i, s_l) \geq u$. If $U$ is uniform on $(0, 1]$, $P(\phi(i, U) = j) = p(i, j)$. Take $X_0 \sim \mu$ and independent uniform
$U_1, U_2, \ldots$ (Corollary~\ref{measure-theory:corollary-independent-sequences}), and $X_{n+1} = \phi(X_n, U_{n+1})$. Since $U_{n+1}$ is independent of $(X_0, \ldots, X_n)$, the defining property holds, and the product formula follows
by induction. Summing it over $i_0, \ldots, i_{n-1}$ gives $\mu P^n$; Chapman--Kolmogorov is the associativity $P^{m+n} = P^mP^n$.
\end{proof}

\begin{example}\label{martingales:example-markov-chains}
\begin{enumerate}
\item \emph{Random walks.} On $\ZZ^d$, $p(x, x \pm e_k) = \frac{1}{2d}$: the \emph{simple random walk}. On a graph with finite degrees, $p(v, w) = 1/\deg v$ for each neighbour $w$ of $v$.
\item \emph{Ehrenfest urn.} $N$ molecules in two chambers; at each step a uniformly chosen molecule changes chamber. The number $X_n$ in the left chamber is a chain on $\{0, \ldots, N\}$ with $p(k, k - 1) = \frac kN$,
$p(k, k + 1) = \frac{N - k}{N}$: a model of the approach to thermal equilibrium.
\item \emph{Card shuffling.} Riffle shuffles or random transpositions of a deck form a random walk on the symmetric group $S_{52}$.
\item \emph{Branching processes} (Definition~\ref{martingales:definition-branching}) and gambler's ruin are Markov chains.
\end{enumerate}
\end{example}

\begin{proposition}[Strong Markov property]\label{martingales:proposition-strong-markov}
Let $T$ be a stopping time of the natural filtration. Conditionally on $\{T < \infty, X_T = i\}$, the process $(X_{T+n})_{n \geq 0}$ is a Markov chain with transition matrix $P$ started at $i$, independent of
$X_0, \ldots, X_T$: for events $A \in \mathcal{F}_T = \{A : A \cap \{T = n\} \in \mathcal{F}_n \text{ for all } n\}$ and $B \subseteq S^{\ZZ_{\geq 0}}$ measurable,
$P_\mu(A, T < \infty, X_T = i, (X_{T+n})_n \in B) = P_\mu(A, T < \infty, X_T = i)P_i((X_n)_n \in B)$.
\end{proposition}

\begin{proof}
Decompose according to $\{T = m\}$: it suffices to prove the identity with $T = m$ fixed and $A \in \mathcal{F}_m$, and for cylinder sets $B = \{x_0 = j_0, \ldots, x_k = j_k\}$, which form a $\pi$-system generating the product
$\sigma$-algebra (Lemma~\ref{measure-theory:lemma-dynkin}). For $A = \{X_0 = i_0, \ldots, X_m = i_m\}$ with $i_m = i$ this is the product formula of Proposition~\ref{martingales:proposition-markov-construction}, and a general
$A \in \mathcal{F}_m$ is a countable disjoint union of such events.
\end{proof}

\begin{definition}\label{martingales:definition-recurrence}
A state $i$ is \emph{recurrent} if $P_i(T_i < \infty) = 1$ and \emph{transient} otherwise; a recurrent state is \emph{positive recurrent} if $E_i[T_i] < \infty$ and \emph{null recurrent} otherwise. State $j$ is
\emph{accessible} from $i$ if $p^n(i, j) > 0$ for some $n \geq 0$; the chain is \emph{irreducible} if every state is accessible from every other. The \emph{period} of $i$ is $\gcd\{n \geq 1 : p^n(i, i) > 0\}$; the chain is
\emph{aperiodic} if every state has period $1$.
\end{definition}

\begin{theorem}\label{martingales:theorem-recurrence-criterion}
A state $i$ is recurrent iff $\sum_{n \geq 1}p^n(i, i) = \infty$. Precisely, the number of returns $N_i = \#\{n \geq 1 : X_n = i\}$ satisfies $P_i(N_i \geq k) = f^k$ with $f = P_i(T_i < \infty)$, and $E_i[N_i] = \sum_{n \geq 1}p^n(i, i)$.
In an irreducible chain all states are recurrent or all are transient, and if they are recurrent, $P_j(T_i < \infty) = 1$ for all $i$, $j$.
\end{theorem}

\begin{proof}
$E_i[N_i] = \sum_nP_i(X_n = i)$ by monotone convergence. By the strong Markov property at the time of the $(k - 1)$-st return, $P_i(N_i \geq k) = fP_i(N_i \geq k - 1)$, so $P_i(N_i \geq k) = f^k$ and $E_i[N_i] = \sum_{k \geq 1}f^k$, which is
finite iff $f < 1$. If $i$ is recurrent and $p^a(i, j) > 0$, $p^b(j, i) > 0$, then $p^{a+n+b}(j, j) \geq p^b(j, i)p^n(i, i)p^a(i, j)$, so $\sum_np^n(j, j) = \infty$ and $j$ is recurrent. Finally, let $i$ be recurrent and $j$ accessible from $i$. Taking a path from $i$ to $j$ of positive probability and cutting it at its last visit to $i$, there is $n$ with $P_i(X_n = j, T_i > n) > 0$.
By the Markov property at time $n$, $0 = P_i(T_i = \infty) \geq P_i(X_n = j, T_i > n)P_j(T_i = \infty)$, so $P_j(T_i < \infty) = 1$.
\end{proof}

\begin{theorem}[P\'olya]\label{martingales:theorem-polya}
The simple random walk on $\ZZ^d$ is recurrent for $d = 1, 2$ and transient for $d \geq 3$.
\end{theorem}

\begin{proof}
Let $\phi(\theta) = \frac1d\sum_{k=1}^d\cos\theta_k$ for $\theta \in Q = [-\pi, \pi]^d$, the characteristic function of a step: $E[e^{i\langle\theta, \xi\rangle}] = \phi(\theta)$. By independence of the steps, $E_0[e^{i\langle\theta, X_n\rangle}] = \phi(\theta)^n$,
and by orthogonality of the characters $e^{i\langle\theta, x\rangle}$, $x \in \ZZ^d$, on $Q$,
\[
p^n(0, 0) = \frac{1}{(2\pi)^d}\int_Q\phi(\theta)^n\,d\theta.
\]
As $p^n(0, 0) = 0$ for odd $n$ and $\phi^{2n} \geq 0$, monotone convergence gives $\sum_np^n(0, 0) = (2\pi)^{-d}\int_Q\frac{d\theta}{1 - \phi(\theta)^2} \in (0, \infty]$. Now $1 - \phi^2 = (1 - \phi)(1 + \phi)$ vanishes on $Q$ only at
$\theta = 0$ (where $\phi = 1$) and at the corners $\theta \in \{\pm\pi\}^d$ (where $\phi = -1$); since $\frac{2t^2}{\pi^2} \leq 1 - \cos t \leq \frac{t^2}{2}$ on $[-\pi, \pi]$, $1 - \phi(\theta)$ is comparable to $|\theta|^2$, and by the symmetry
$\phi(\theta + (\pi, \ldots, \pi)) = -\phi(\theta)$ the same holds for $1 + \phi$ near the corners. Away from these points $1 - \phi^2$ is bounded below. So the integral is finite iff $\int_{|\theta| < 1}|\theta|^{-2}\,d\theta < \infty$, that is,
iff $d \geq 3$ (polar coordinates). Theorem~\ref{martingales:theorem-recurrence-criterion} concludes.
\end{proof}

\begin{remark}\label{martingales:remark-polya}
``A drunk man will find his way home, but a drunk bird may get lost forever'' (Kakutani). In $\ZZ^3$ the return probability is $1 - 1/u_3 \approx 0.3405$, where $u_3 = (2\pi)^{-3}\int_Q\frac{d\theta}{1 - \phi(\theta)}$ is the expected number of
visits to $0$ (Watson's integral). In $\ZZ^1$ the walk is null recurrent: $E_0[T_0] = \infty$, by Counterexample~\ref{martingales:counterexample-optional-stopping}.
\end{remark}

\begin{definition}\label{martingales:definition-stationary}
A \emph{stationary distribution} is a probability measure $\pi$ on $S$ with $\pi P = \pi$. A measure $\pi$ is \emph{reversible} if $\pi(i)p(i, j) = \pi(j)p(j, i)$ for all $i$, $j$ (\emph{detailed balance}).
\end{definition}

A reversible probability measure is stationary: $\sum_i\pi(i)p(i, j) = \sum_i\pi(j)p(j, i) = \pi(j)$. If $X_0 \sim \pi$ is stationary, then $X_n \sim \pi$ for all $n$, and the chain is a stationary process; if $\pi$ is reversible, the
chain run backwards in time has the same law.

\begin{example}\label{martingales:example-stationary}
\begin{enumerate}
\item For the random walk on a finite connected graph $G = (V, E)$, $\pi(v) = \frac{\deg v}{2|E|}$ is reversible: $\pi(v)p(v, w) = \frac{1}{2|E|}$ for adjacent $v$, $w$.
\item For the Ehrenfest urn, the binomial law $\pi(k) = \binom Nk2^{-N}$ is reversible: $\binom Nk\frac{N - k}{N} = \binom N{k+1}\frac{k + 1}{N}$.
\item The simple random walk on $\ZZ^d$ has no stationary distribution: the invariant measures are multiples of counting measure, which is infinite.
\end{enumerate}
\end{example}

\begin{theorem}\label{martingales:theorem-stationary}
Let the chain be irreducible and recurrent, and fix a state $k$. Then $\lambda(j) = E_k\sum_{n=0}^{T_k - 1}\mathbf{1}_{\{X_n = j\}}$, the expected time spent in $j$ during an excursion from $k$, defines an invariant measure ($\lambda P = \lambda$) with
$0 < \lambda(j) < \infty$ and $\lambda(k) = 1$, and every invariant measure is a multiple of $\lambda$. Consequently, an irreducible chain has a stationary distribution iff it is positive recurrent, and then the stationary distribution
is unique, given by $\pi(j) = 1/E_j[T_j]$.
\end{theorem}

\begin{proof}[Sketch of proof]
\emph{Invariance.} $\lambda(j) = \sum_{n \geq 1}P_k(X_n = j, n \leq T_k)$ (shifting the excursion by one step, using $X_0 = X_{T_k} = k$), and by the Markov property at time $n - 1$, with $\{n \leq T_k\} \in \mathcal{F}_{n-1}$,
$P_k(X_n = j, n \leq T_k) = \sum_iP_k(X_{n-1} = i, n - 1 < T_k)p(i, j)$; summing over $n$ gives $\lambda(j) = \sum_i\lambda(i)p(i, j)$. \emph{Finiteness and positivity.} From $\lambda = \lambda P^n$, $1 = \lambda(k) \geq \lambda(j)p^n(j, k)$ and
$\lambda(j) \geq \lambda(k)p^n(k, j)$, for suitable $n$. \emph{Uniqueness.} If $\mu$ is invariant with $\mu(k) = 1$, iterating $\mu(j) = \mu(k)p(k, j) + \sum_{i \neq k}\mu(i)p(i, j)$ shows $\mu \geq \lambda$, and $\mu - \lambda$ is invariant and
vanishes at $k$, hence everywhere by the positivity argument. \emph{Stationary distributions.} $\sum_j\lambda(j) = E_k[T_k]$, so $\lambda$ normalises iff $k$ is positive recurrent, and then $\pi(k) = \lambda(k)/E_k[T_k] = 1/E_k[T_k]$. If
the chain is transient, $p^n(i, j) \to 0$ for all $i$, $j$ (the series $\sum_np^n(i, j) = P_i(T_j < \infty)E_j[N_j + 1]$ converges), so a stationary $\pi = \pi P^n$ would vanish by dominated convergence. The omitted verifications are
the interchanges of sums, justified by nonnegativity.
\end{proof}

\begin{example}\label{martingales:example-knight}
A knight moves at random on a chessboard, choosing each legal move with equal probability. Starting from a corner, the expected number of moves until it returns is $1/\pi(\text{corner}) = 2|E|/\deg(\text{corner}) = 336/2 = 168$,
by Example~\ref{martingales:example-stationary}(1) and Exercise~\ref{martingales:exercise-knight}.
\end{example}

\begin{lemma}\label{martingales:lemma-aperiodic}
If $i$ has period $1$, then $p^n(i, i) > 0$ for all large $n$.
\end{lemma}

\begin{proof}
$A = \{n \geq 1 : p^n(i, i) > 0\}$ is closed under addition (Chapman--Kolmogorov). Its gcd is $1$, so some finite subset has gcd $1$, and by B\'ezout's theorem (Theorem~\ref{number-systems:theorem-bezout}, applied repeatedly) there are integers $c_a$ with
$\sum_ac_aa = 1$. Collecting the terms with $c_a > 0$ and those with $c_a < 0$, $1 = x - y$ where $x$ and $y$ are sums of elements of $A$ (or $0$), hence in $A \cup \{0\}$. If $y = 0$, then $1 \in A$ and $A = \ZZ_{\geq 1}$. Otherwise
$x = y + 1$ with $x, y \in A$. For $n \geq y^2$ write $n = qy + r$ with $0 \leq r < y \leq q$; then $n = (q - r)y + rx$ with $q - r \geq 1$, so $n \in A$.
\end{proof}

\begin{theorem}[Convergence to equilibrium]\label{martingales:theorem-convergence-equilibrium}
Let the chain be irreducible, aperiodic, with stationary distribution $\pi$. Then for every initial distribution $\mu$, $d_{TV}(\mu P^n, \pi) \to 0$; in particular $p^n(i, j) \to \pi(j)$.
\end{theorem}

\begin{proof}
\emph{Coupling.} Let $(X_n)$ and $(Y_n)$ be independent chains with $X_0 \sim \mu$ and $Y_0 \sim \pi$. The pair $(X_n, Y_n)$ is a Markov chain on $S \times S$ with transition probabilities $p(i, j)p(k, l)$ and stationary distribution
$\pi \otimes \pi$. It is irreducible: by Lemma~\ref{martingales:lemma-aperiodic} and irreducibility, $p^n(i, j)p^n(k, l) > 0$ for all large $n$. By Theorem~\ref{martingales:theorem-stationary} it is positive recurrent, so by
Theorem~\ref{martingales:theorem-recurrence-criterion} the time $\tau = \inf\{n : X_n = Y_n\}$ is almost surely finite.
\emph{Swapping.} Let $Z_n = X_n$ for $n < \tau$ and $Z_n = Y_n$ for $n \geq \tau$. By the strong Markov property at $\tau$, $(Z_n)$ is again a chain with transition matrix $P$ and $Z_0 \sim \mu$, so $Z_n \sim \mu P^n$, while $Y_n \sim \pi$.
Hence for every $A \subseteq S$, $|\mu P^n(A) - \pi(A)| = |P(Z_n \in A) - P(Y_n \in A)| \leq P(Z_n \neq Y_n) \leq P(\tau > n) \to 0$.
\end{proof}

\begin{counterexample}\label{martingales:counterexample-periodic}
Aperiodicity is needed: the walk on two states with $p(1, 2) = p(2, 1) = 1$ has stationary distribution $(\frac12, \frac12)$, but $p^n(1, 1)$ alternates between $0$ and $1$. The \emph{lazy} chain $\frac12(I + P)$ is aperiodic with
the same stationary distribution, which is the standard cure.
\end{counterexample}

\begin{theorem}[Ergodic theorem for Markov chains]\label{martingales:theorem-markov-ergodic}
Let the chain be irreducible and positive recurrent with stationary distribution $\pi$, and $f \colon S \to \RR$ with $\sum_j\pi(j)|f(j)| < \infty$. Then for every initial distribution,
$\frac1n\sum_{k=0}^{n-1}f(X_k) \to \sum_j\pi(j)f(j)$ almost surely. In particular the fraction of time spent in $j$ tends to $\pi(j)$.
\end{theorem}

\begin{proof}[Sketch of proof]
Let $\tau_0 < \tau_1 < \cdots$ be the successive visits to a fixed state $k$ ($\tau_0 < \infty$ a.s.\ by Theorem~\ref{martingales:theorem-recurrence-criterion}). By the strong Markov property the excursions between them are
i.i.d., so the sums $Y_r = \sum_{\tau_{r-1} \leq n < \tau_r}f(X_n)$ and the lengths $\tau_r - \tau_{r-1}$ are i.i.d.\ sequences with means $\sum_j\lambda(j)f(j)$ and $E_k[T_k]$, in the notation of Theorem~\ref{martingales:theorem-stationary}. The
strong law (Theorem~\ref{probability:theorem-strong-law}) applied to both, together with the bound of the partial excursion at time $n$ (handled by applying the same argument to $|f|$), gives the limit
$\sum_j\lambda(j)f(j)/E_k[T_k] = \sum_j\pi(j)f(j)$.
\end{proof}

\begin{remark}\label{martingales:remark-mixing}
\begin{enumerate}
\item \emph{Markov chain Monte Carlo.} To sample from a probability measure $\pi$ on a huge finite set (Ising configurations, colourings of a graph), one designs a chain with stationary distribution $\pi$ and runs it for a
long time. The \emph{Metropolis algorithm}: propose a move $i \to j$ by a symmetric matrix $Q$ and accept it with probability $\min(1, \pi(j)/\pi(i))$; the resulting chain is reversible with respect to $\pi$, and only ratios of
$\pi$ are needed, not its normalising constant.
\item \emph{Mixing times.} The \emph{mixing time} $t_{\mathrm{mix}} = \min\{n : \max_id_{TV}(p^n(i, \cdot), \pi) \leq \frac14\}$ measures the speed of Theorem~\ref{martingales:theorem-convergence-equilibrium}. For a reversible chain the
eigenvalues of $P$ are real, $1 = \lambda_1 > \lambda_2 \geq \cdots \geq \lambda_{|S|} \geq -1$, and $t_{\mathrm{mix}}$ is controlled by the \emph{spectral gap} $1 - \max(\lambda_2, |\lambda_{|S|}|)$. Bayer and Diaconis showed that
$\frac32\log_2n$ riffle shuffles mix a deck of $n$ cards, with a sharp \emph{cutoff}; for $52$ cards, seven shuffles suffice \cite{levin-peres-wilmer}.
\item \emph{Electrical networks.} For the random walk on a graph with edge conductances, hitting probabilities are voltages and escape probabilities are effective conductances; recurrence is equivalent to infinite
effective resistance to infinity, which gives a robust proof of P\'olya's theorem (Rayleigh's monotonicity) \cite{doyle-snell}.
\end{enumerate}
\end{remark}

\section{The Poisson process}\label{martingales:section-poisson-process}

\begin{definition}\label{martingales:definition-poisson-process}
Let $\tau_1, \tau_2, \ldots$ be i.i.d.\ $\mathrm{Exp}(\lambda)$ and $T_n = \tau_1 + \cdots + \tau_n$. The \emph{Poisson process} of rate $\lambda$ is $N_t = \#\{n \geq 1 : T_n \leq t\}$, $t \geq 0$: it counts the arrivals up to time $t$
when the waiting times between arrivals are independent exponentials.
\end{definition}

\begin{theorem}\label{martingales:theorem-poisson-process}
The Poisson process has independent stationary Poisson increments: for $0 = t_0 < t_1 < \cdots < t_k$ the increments $N_{t_j} - N_{t_{j-1}}$ are independent with laws $\mathrm{Poi}(\lambda(t_j - t_{j-1}))$. Given $N_t = n$, the
arrival times $T_1, \ldots, T_n$ are distributed as the order statistics of $n$ independent uniform points of $[0, t]$.
\end{theorem}

\begin{proof}
\emph{Joint density.} $(T_1, \ldots, T_{n+1})$ has density $\lambda^{n+1}e^{-\lambda s_{n+1}}$ on $\{0 < s_1 < \cdots < s_{n+1}\}$, by the change of variables $s_j = \tau_1 + \cdots + \tau_j$ (Jacobian $1$). \emph{Conditional law.} For a
Borel set $A \subseteq \{0 < s_1 < \cdots < s_n < t\}$,
\[
P((T_1, \ldots, T_n) \in A, N_t = n) = \int_A\int_t^\infty\lambda^{n+1}e^{-\lambda s_{n+1}}\,ds_{n+1}\,ds = \lambda^ne^{-\lambda t}|A|.
\]
With $A$ the whole simplex, of volume $t^n/n!$, this gives $P(N_t = n) = e^{-\lambda t}(\lambda t)^n/n!$, and dividing, the conditional density of $(T_1, \ldots, T_n)$ is the constant $n!/t^n$, which is the density of the order statistics
of $n$ uniform points. \emph{Increments.} Given $N_t = n$, the numbers of the $n$ uniform points falling in the intervals $(t_{j-1}, t_j]$, $t_k = t$, are multinomial with probabilities $(t_j - t_{j-1})/t$. Multiplying the multinomial
probabilities by $P(N_t = n)$ factorises as a product of Poisson probabilities, as in Exercise~\ref{probability:exercise-poisson-thinning}.
\end{proof}

\begin{remark}\label{martingales:remark-poisson-process}
\begin{enumerate}
\item The Poisson process is the simplest continuous-time Markov chain: it jumps from $n$ to $n + 1$ at rate $\lambda$. A general continuous-time chain on a countable set waits an exponential time with a rate $q(i)$ in state $i$ and
then jumps according to a transition matrix; it is described by its \emph{generator} $Q$, and its transition probabilities are $e^{tQ}$.
\item $N_t - \lambda t$ and $(N_t - \lambda t)^2 - \lambda t$ are martingales; by the law of rare events (Theorem~\ref{probability:theorem-poisson-limit}), the Poisson process is the limit of Bernoulli trials with success probability
$\lambda\,\Delta t$ in each small time interval. Superposition of independent Poisson processes and independent thinning give Poisson processes again (Exercise~\ref{probability:exercise-poisson-thinning}).
\item The arrival times form a \emph{Poisson point process} on $[0, \infty)$; more generally, for a measure $\mu$ on a space $E$ there is a random point set whose counts in disjoint sets are independent with Poisson laws of means
$\mu(A)$. This is the jump structure of L\'evy processes (Theorem~\ref{probability:theorem-levy-khintchine}).
\end{enumerate}
\end{remark}

\section{Brownian motion}\label{martingales:section-brownian-motion}

\begin{definition}\label{martingales:definition-brownian-motion}
A (standard, one-dimensional) \emph{Brownian motion} is a process $(B_t)_{t \geq 0}$ on a probability space such that
\begin{enumerate}
\item $B_0 = 0$;
\item for $0 \leq t_0 < t_1 < \cdots < t_k$ the increments $B_{t_j} - B_{t_{j-1}}$ are independent, and $B_t - B_s \sim N(0, t - s)$ for $s < t$;
\item for every $\omega$, $t \mapsto B_t(\omega)$ is continuous.
\end{enumerate}
A $d$-dimensional Brownian motion is $(B^1_t, \ldots, B^d_t)$ with independent one-dimensional Brownian motions $B^k$.
\end{definition}

By the central limit theorem, (2) is forced on any process with independent stationary increments of finite variance and continuous paths, and Brownian motion is the scaling limit of random walks
(Remark~\ref{martingales:remark-donsker}). It was used by Bachelier (1900) for stock prices and by Einstein (1905) for the motion of pollen particles, and constructed rigorously by Wiener (1923).

\begin{theorem}[Wiener; L\'evy's construction]\label{martingales:theorem-brownian-existence}
Brownian motion exists.
\end{theorem}

\begin{proof}[Sketch of proof]
It suffices to construct it on $[0, 1]$ and to concatenate independent copies. Let $Z_0$ and $Z_{n,k}$ ($n \geq 0$, $k$ odd, $0 < k < 2^n$) be independent $N(0, 1)$ variables
(Corollary~\ref{measure-theory:corollary-independent-sequences}). Define $B$ at the dyadic points by induction: $B_0 = 0$, $B_1 = Z_0$, and at a new dyadic point $t = k2^{-n}$ ($k$ odd), with neighbours $t_\pm = t \pm 2^{-n}$,
\[
B_t = \tfrac12(B_{t_-} + B_{t_+}) + 2^{-(n+1)/2}Z_{n,k}.
\]
\emph{Increments.} Induction on $n$ shows that the increments of $B$ between consecutive points of $2^{-n}\ZZ \cap [0, 1]$ are independent $N(0, 2^{-n})$: if $X = B_{t_+} - B_{t_-} \sim N(0, 2^{1-n})$ and $Y = 2^{-(n+1)/2}Z_{n,k}$, then
$\frac X2 + Y$ and $\frac X2 - Y$ are independent $N(0, 2^{-n})$ variables (a computation with characteristic functions in two variables, which is the verification we omit). \emph{Continuity.} The piecewise linear interpolations
$B^{(n)}$ satisfy $\|B^{(n)} - B^{(n-1)}\|_\infty = 2^{-(n+1)/2}\max_k|Z_{n,k}|$, and $P(\max_k|Z_{n,k}| > n) \leq 2^ne^{-n^2/2}$ is summable; by Borel--Cantelli the $B^{(n)}$ converge uniformly to a continuous limit, which extends $B$ from
the dyadic points. \emph{Extension.} Increments over arbitrary times are limits of dyadic increments, and limits in probability of independent Gaussian vectors are Gaussian with the limiting covariance (by L\'evy's continuity
theorem in several variables).
\end{proof}

\begin{proposition}[Invariances]\label{martingales:proposition-brownian-invariance}
If $B$ is a Brownian motion, so are $-B_t$; $c^{-1/2}B_{ct}$ for $c > 0$ (\emph{scaling}); $B_{s+t} - B_s$ for fixed $s \geq 0$, which is moreover independent of $(B_r)_{r \leq s}$ (\emph{simple Markov property}).
\end{proposition}

\begin{proof}
Each process has continuous paths starting at $0$, and its increments over disjoint intervals are increments of $B$ over disjoint intervals, multiplied by constants; their laws are computed from
Example~\ref{probability:example-characteristic-functions}(2): $c^{-1/2}N(0, c(t - s)) = N(0, t - s)$. Independence of $(B_{s+t} - B_s)_t$ from $(B_r)_{r \leq s}$ holds for finitely many times by (2), hence for the generated
$\sigma$-algebras by Lemma~\ref{probability:lemma-pi-system-independence}.
\end{proof}

\begin{proposition}[Quadratic variation]\label{martingales:proposition-quadratic-variation}
Let $0 = t_0 < t_1 < \cdots < t_m = t$ be a partition with mesh $\delta = \max(t_k - t_{k-1})$. Then
\[
E\Bigl[\Bigl(\sum_k(B_{t_k} - B_{t_{k-1}})^2 - t\Bigr)^2\Bigr] = 2\sum_k(t_k - t_{k-1})^2 \leq 2t\delta.
\]
So the sums of squared increments converge to $t$ in $L^2$ as the mesh tends to $0$, and almost surely along any sequence of partitions with summable meshes.
\end{proposition}

\begin{proof}
The terms $D_k = (B_{t_k} - B_{t_{k-1}})^2 - (t_k - t_{k-1})$ are independent with mean $0$ and variance $\operatorname{Var}(\Delta_k^2) = 2(\Delta t_k)^2$, as $E[Z^4] = 3$ for $Z \sim N(0, 1)$. So $E[(\sum D_k)^2] = \sum E[D_k^2]$. For the almost sure
statement, Chebyshev's inequality and Borel--Cantelli.
\end{proof}

\begin{corollary}\label{martingales:corollary-unbounded-variation}
Almost surely, the Brownian path has infinite variation on every interval $[s, t]$, $s < t$.
\end{corollary}

\begin{proof}
It suffices to treat rational $s < t$, a countable family. For dyadic partitions of $[s, t]$ of mesh $2^{-n}(t - s)$, $\sum_k(\Delta_kB)^2 \to t - s > 0$ almost surely by Proposition~\ref{martingales:proposition-quadratic-variation} (applied to the
Brownian motion $B_{s+r} - B_s$). But $\sum_k(\Delta_kB)^2 \leq \max_k|\Delta_kB|\cdot\sum_k|\Delta_kB|$, and $\max_k|\Delta_kB| \to 0$ by uniform continuity of the path on $[s, t]$; so $\sum_k|\Delta_kB| \to \infty$.
\end{proof}

This is why the Riemann--Stieltjes integral $\int f\,dB$ cannot be defined pathwise, and a new integral is needed (Section~\ref{martingales:section-stochastic-calculus}).

\begin{theorem}[Paley--Wiener--Zygmund]\label{martingales:theorem-nowhere-differentiable}
Almost surely, $t \mapsto B_t$ is differentiable at no point of $[0, \infty)$.
\end{theorem}

\begin{proof}
It suffices to consider $[0, 1]$. If the path is differentiable at $s \in [0, 1]$, there are integers $C$, $M$ with $|B_t - B_s| \leq C|t - s|$ for $|t - s| \leq 4/M$. For $n \geq M$, let $j = \lfloor ns\rfloor$ and $\Delta_i = B_{(i+1)/n} - B_{i/n}$;
the points $\frac in$ for $j + 1 \leq i \leq j + 4$ lie within $\frac4n$ of $s$, so $|\Delta_i| \leq \frac{8C}{n}$ for $i = j + 1, j + 2, j + 3$. So the path lies in
$E_n = \{\exists j \leq n : |\Delta_{j+1}|, |\Delta_{j+2}|, |\Delta_{j+3}| \leq 8C/n\}$ for all $n \geq M$. The $\Delta_i$ are independent $N(0, \frac1n)$, whose density is at most $\sqrt{n/2\pi}$, so
\[
P(E_n) \leq (n + 1)\Bigl(\frac{16C}{n}\sqrt{\frac{n}{2\pi}}\Bigr)^3 = O(n^{-1/2}).
\]
Hence $P(\bigcap_{n \geq M}E_n) = 0$, and the set of paths differentiable somewhere is contained in the countable union over $C$, $M$ of these null events.
\end{proof}

\begin{theorem}[Reflection principle]\label{martingales:theorem-reflection}
For $a > 0$ and $t > 0$, $P(\max_{s \leq t}B_s \geq a) = 2P(B_t \geq a) = P(|B_t| \geq a)$. Consequently the hitting time $T_a = \inf\{s : B_s = a\}$ is almost surely finite, with density $\frac{a}{\sqrt{2\pi s^3}}e^{-a^2/2s}$ and
$E[T_a] = \infty$.
\end{theorem}

\begin{proof}[Sketch of proof]
$T_a$ is a stopping time, and by the \emph{strong Markov property} of Brownian motion (the analogue of Proposition~\ref{martingales:proposition-strong-markov}: $(B_{T_a + s} - a)_{s \geq 0}$ is a Brownian motion independent of the past
up to $T_a$, which we do not prove here \cite{morters-peres}), the path reflected at $a$ after $T_a$ is again a Brownian motion. Paths with $\max_{s \leq t}B_s \geq a$ and $B_t < a$ correspond bijectively to paths with $B_t > a$ under
this reflection, so $P(\max \geq a, B_t < a) = P(B_t > a)$; and $P(B_t = a) = 0$. The density of $T_a$ is the derivative of $P(T_a \leq s) = P(|B_s| \geq a) = 2\int_{a/\sqrt s}^\infty\frac{e^{-x^2/2}}{\sqrt{2\pi}}dx$, which tends to $1$ as
$s \to \infty$; and $\int^\infty s\cdot s^{-3/2}ds = \infty$.
\end{proof}

\begin{remark}\label{martingales:remark-brownian}
\begin{enumerate}
\item \emph{Martingales.} $B_t$, $B_t^2 - t$ and $\exp(\lambda B_t - \lambda^2t/2)$ are continuous-time martingales. Optional stopping as in Example~\ref{martingales:example-gamblers-ruin} shows that Brownian motion started at $0$
leaves $(-b, a)$ at $a$ with probability $\frac{b}{a + b}$, after an expected time $ab$.
\item \emph{Harmonic functions.} For a bounded domain $D \subseteq \RR^d$ with regular boundary and continuous $g$ on $\partial D$, $u(x) = E_x[g(B_\tau)]$, where $\tau$ is the exit time of $D$, is the solution of the Dirichlet problem
$\Delta u = 0$ in $D$, $u = g$ on $\partial D$ (Kakutani, 1944); the mean value property of harmonic functions is the rotation invariance of Brownian motion. Recurrence of Brownian motion in $\RR^2$ and transience in $\RR^d$,
$d \geq 3$, mirror P\'olya's theorem, and correspond to the harmonic functions $\log|x|$ and $|x|^{2-d}$.
\item \emph{Law of the iterated logarithm.} $\limsup_{t \to \infty}B_t/\sqrt{2t\log\log t} = 1$ almost surely, and by time inversion ($tB_{1/t}$ is a Brownian motion) $\limsup_{t \to 0}B_t/\sqrt{2t\log\log(1/t)} = 1$: paths are H\"older
continuous of every order $<\frac12$ but not of order $\frac12$ \cite{morters-peres}.
\end{enumerate}
\end{remark}

\begin{remark}[Donsker's invariance principle]\label{martingales:remark-donsker}
If $S_n$ is a random walk with i.i.d.\ steps of mean $0$ and variance $1$, the rescaled interpolated paths $t \mapsto S_{\lfloor nt\rfloor}/\sqrt n$ converge in distribution, as random elements of the space of paths with the uniform
topology, to Brownian motion. Hence functionals of random walks, such as $\max_{k \leq n}S_k/\sqrt n$ or the fraction of time spent positive (Paul L\'evy's arcsine law), converge in distribution to the corresponding Brownian
functionals; tightness is proved with the Arzel\`a--Ascoli theorem (Theorem~\ref{general-topology:theorem-arzela-ascoli}) and Prokhorov's theorem \cite{morters-peres}.
\end{remark}

\section{Stochastic calculus}\label{martingales:section-stochastic-calculus}

Let $B$ be a Brownian motion and $\mathcal{F}_t = \sigma(B_s : s \leq t)$. We want to define $\int_0^tH_s\,dB_s$ for random integrands $H$ that do not anticipate the future, as the continuous-time analogue of the martingale
transform of Definition~\ref{martingales:definition-transform}.

\begin{definition}\label{martingales:definition-simple-process}
A \emph{simple process} is $H_s = \sum_{k=0}^{m-1}\eta_k\mathbf{1}_{(t_k, t_{k+1}]}(s)$ with $0 = t_0 < \cdots < t_m$ and bounded $\mathcal{F}_{t_k}$-measurable random variables $\eta_k$. Its \emph{It\^o integral} is
$\int_0^\infty H_s\,dB_s = \sum_k\eta_k(B_{t_{k+1}} - B_{t_k})$.
\end{definition}

\begin{proposition}[It\^o isometry]\label{martingales:proposition-ito-isometry}
For simple $H$, $E[\int H\,dB] = 0$ and $E\bigl[\bigl(\int_0^\infty H_s\,dB_s\bigr)^2\bigr] = E\int_0^\infty H_s^2\,ds$.
\end{proposition}

\begin{proof}
Let $\Delta_k = B_{t_{k+1}} - B_{t_k}$, independent of $\mathcal{F}_{t_k}$ with mean $0$ and variance $t_{k+1} - t_k$. Then $E[\eta_k\Delta_k] = E[\eta_k]E[\Delta_k] = 0$. For $j < k$, $\eta_j\Delta_j\eta_k$ is $\mathcal{F}_{t_k}$-measurable, so
$E[\eta_j\Delta_j\eta_k\Delta_k] = 0$; and $E[\eta_k^2\Delta_k^2] = E[\eta_k^2](t_{k+1} - t_k)$. Summing, $E[(\sum_k\eta_k\Delta_k)^2] = \sum_kE[\eta_k^2](t_{k+1} - t_k) = E\int H^2\,ds$.
\end{proof}

\begin{theorem}\label{martingales:theorem-ito-integral}
Let $\mathcal{L}^2$ be the space of processes $H$, measurable in $(s, \omega)$ and adapted ($H_s$ is $\mathcal{F}_s$-measurable), with $E\int_0^\infty H_s^2\,ds < \infty$. Simple processes are dense in $\mathcal{L}^2$ for the norm
$(E\int H^2\,ds)^{1/2}$, so the It\^o integral extends uniquely to an isometry $\mathcal{L}^2 \to L^2(\Omega)$. The process $t \mapsto \int_0^tH_s\,dB_s = \int H_s\mathbf{1}_{[0, t]}(s)\,dB_s$ has a continuous version, which is a martingale.
\end{theorem}

\begin{proof}[Sketch of proof]
Given the density, the extension is the extension of a uniformly continuous linear map from a dense subspace into a complete space. For density: bounded $H$ with continuous paths are approximated by the left-endpoint simple
processes $\sum_kH_{t_k}\mathbf{1}_{(t_k, t_{k+1}]}$, by dominated convergence; general $H$ are approximated by such processes by truncation and convolution in time with a one-sided mollifier, which preserves adaptedness.
The martingale property holds for simple $H$ by Proposition~\ref{martingales:proposition-conditional-expectation-properties} and passes to $L^2$ limits; continuity follows from Doob's $L^2$ inequality
(Theorem~\ref{martingales:theorem-doob-lp}, in its continuous-time form) and Borel--Cantelli. These are the verifications we omit \cite{karatzas-shreve}.
\end{proof}

\begin{example}\label{martingales:example-ito-bdb}
$\int_0^tB_s\,dB_s = \frac12(B_t^2 - t)$. Indeed, for a partition $0 = t_0 < \cdots < t_m = t$, the left-endpoint simple process $H^{(m)} = \sum_kB_{t_k}\mathbf{1}_{(t_k, t_{k+1}]}$ satisfies
$E\int_0^t(B_s - H^{(m)}_s)^2ds = \sum_k\frac12(t_{k+1} - t_k)^2 \to 0$, so $\int H^{(m)}\,dB \to \int_0^tB\,dB$ in $L^2$; and
\[
\int H^{(m)}\,dB = \sum_kB_{t_k}(B_{t_{k+1}} - B_{t_k}) = \frac12B_t^2 - \frac12\sum_k(B_{t_{k+1}} - B_{t_k})^2 \to \frac12(B_t^2 - t)
\]
in $L^2$ by Proposition~\ref{martingales:proposition-quadratic-variation}. The extra term $-\frac t2$, absent from the ordinary calculus $\int x\,dx = \frac{x^2}{2}$, comes from the quadratic variation; with midpoint sums one would get
the \emph{Stratonovich} integral $\frac12B_t^2$, which obeys the ordinary chain rule but is not a martingale.
\end{example}

\begin{theorem}[It\^o's formula]\label{martingales:theorem-ito-formula}
For $f \in C^2(\RR)$ with bounded $f^{\prime}$ and $f^{\prime\prime}$,
\[
f(B_t) = f(B_0) + \int_0^tf^{\prime}(B_s)\,dB_s + \frac12\int_0^tf^{\prime\prime}(B_s)\,ds.
\]
More generally, for $f \in C^{1,2}([0, \infty) \times \RR)$ with bounded derivatives, $df(t, B_t) = \partial_tf\,dt + \partial_xf\,dB_t + \frac12\partial_x^2f\,dt$.
\end{theorem}

\begin{proof}[Sketch of proof]
For a partition of $[0, t]$, Taylor's formula gives $f(B_t) - f(B_0) = \sum_kf^{\prime}(B_{t_k})\Delta_kB + \frac12\sum_kf^{\prime\prime}(B_{t_k})(\Delta_kB)^2 + R$, where $|R| \leq \sum_k\omega(\max_j|\Delta_jB|)(\Delta_kB)^2$ with $\omega$ the modulus
of continuity of $f^{\prime\prime}$ on the range of the path. As the mesh tends to $0$: the first sum tends to $\int f^{\prime}(B)\,dB$ in $L^2$, as in Example~\ref{martingales:example-ito-bdb}; the second tends to $\frac12\int f^{\prime\prime}(B_s)\,ds$,
by the argument of Proposition~\ref{martingales:proposition-quadratic-variation} applied with the weights $f^{\prime\prime}(B_{t_k})$; and $R \to 0$ in probability. The omitted verification is the weighted quadratic variation estimate
\cite{karatzas-shreve}.
\end{proof}

Informally: $(dB_t)^2 = dt$, $dt\,dB_t = 0$, $(dt)^2 = 0$, and Taylor's formula is kept to second order. With $f(x) = x^2$ the formula recovers Example~\ref{martingales:example-ito-bdb}.

\begin{example}[Geometric Brownian motion]\label{martingales:example-geometric-brownian}
$X_t = X_0\exp(\sigma B_t + (\mu - \frac12\sigma^2)t)$ satisfies, by It\^o's formula with $f(t, x) = X_0e^{\sigma x + (\mu - \sigma^2/2)t}$,
\[
dX_t = \mu X_t\,dt + \sigma X_t\,dB_t,
\]
the \emph{stochastic differential equation} of a price with constant expected return $\mu$ and volatility $\sigma$. Note $E[X_t] = X_0e^{\mu t}$, while the typical growth rate is $\mu - \frac{\sigma^2}{2}$: $\frac1t\log X_t \to \mu - \frac{\sigma^2}{2}$
almost surely, by the law of large numbers for $B_t/t$.
\end{example}

\begin{remark}\label{martingales:remark-stochastic-calculus}
\begin{enumerate}
\item \emph{Stochastic differential equations.} $dX_t = b(X_t)\,dt + \sigma(X_t)\,dB_t$ with Lipschitz $b$, $\sigma$ has a unique strong solution, by Picard iteration in the space of continuous adapted processes with the It\^o isometry and
Doob's inequality replacing the sup norm estimates of the deterministic case (Theorem~\ref{real-analysis:theorem-picard-lindelof}). The solution is a Markov process with generator $Lf = bf^{\prime} + \frac12\sigma^2f^{\prime\prime}$; the
\emph{Ornstein--Uhlenbeck process} $dX = -\theta X\,dt + \sigma\,dB$ is the Gaussian stationary example \cite{oksendal}.
\item \emph{Girsanov's theorem.} Under the probability $dQ = \exp(\theta B_T - \frac12\theta^2T)\,dP$, the process $B_t - \theta t$, $t \leq T$, is a Brownian motion: drift can be removed by a change of measure. \emph{L\'evy's
characterisation}: a continuous martingale $M$ with $M_0 = 0$ and $M_t^2 - t$ a martingale is a Brownian motion. \emph{Martingale representation}: every square integrable $\mathcal{F}_T$-measurable variable is
$E[X] + \int_0^TH\,dB$ \cite{revuz-yor}.
\item \emph{Feynman--Kac.} The solution of $\partial_tu = \frac12\Delta u - Vu$, $u(0, \cdot) = g$, is $u(t, x) = E_x[g(B_t)e^{-\int_0^tV(B_s)\,ds}]$: parabolic PDE and Brownian expectations are the same thing.
\item \emph{Black--Scholes.} If a share follows geometric Brownian motion and money earns interest at rate $r$, the arbitrage-free price of a European call option with strike $K$ and maturity $T$ is the expectation of its
discounted payoff under the measure making $e^{-rt}X_t$ a martingale (Girsanov); this gives the Black--Scholes formula
$C = X_0\Phi(d_1) - Ke^{-rT}\Phi(d_2)$ with $d_{1,2} = \frac{\log(X_0/K) + (r \pm \sigma^2/2)T}{\sigma\sqrt T}$ \cite{black-scholes-1973}.
\end{enumerate}
\end{remark}

\section{Exercises}\label{martingales:section-exercises}

\begin{exercise}\label{martingales:exercise-projection}
Let $X \in L^2$ and $\mathcal{G}$ a sub-$\sigma$-algebra. Show that $E[(X - Y)^2] \geq E[(X - E[X \mid \mathcal{G}])^2]$ for every $\mathcal{G}$-measurable $Y \in L^2$, and that $\operatorname{Var}X = E[\operatorname{Var}(X \mid \mathcal{G})] + \operatorname{Var}E[X \mid \mathcal{G}]$, where
$\operatorname{Var}(X \mid \mathcal{G}) = E[X^2 \mid \mathcal{G}] - E[X \mid \mathcal{G}]^2$.
\end{exercise}

\begin{solution}
With $Z = E[X \mid \mathcal{G}]$, $X - Z \perp L^2(\mathcal{G})$ (Theorem~\ref{martingales:theorem-conditional-expectation}), so $E[(X - Y)^2] = E[(X - Z)^2] + E[(Z - Y)^2]$ by Pythagoras. For the variance formula,
$E[\operatorname{Var}(X \mid \mathcal{G})] = E[X^2] - E[Z^2]$ and $\operatorname{Var}Z = E[Z^2] - E[X]^2$ (using $E[Z] = E[X]$); add.
\end{solution}

\begin{exercise}\label{martingales:exercise-hth}
A fair coin is tossed until the pattern $HTH$ appears. Show that the expected number of tosses is $10$, whereas for $HHT$ it is $8$.
\end{exercise}

\begin{solution}
Let $e_0$, $e_1$, $e_2$ be the expected remaining numbers of tosses when the current progress towards $HTH$ is nothing, $H$, $HT$. Conditioning on the next toss (a Markov chain argument): $e_0 = 1 + \frac12e_1 + \frac12e_0$,
$e_1 = 1 + \frac12e_1 + \frac12e_2$, $e_2 = 1 + \frac12\cdot 0 + \frac12e_0$. So $e_0 = 2 + e_1$, $e_1 = 2 + e_2$, $e_2 = 1 + \frac12e_0$, whence $e_0 = 5 + \frac12e_0$ and $e_0 = 10$. For $HHT$: $e_0 = 1 + \frac12e_1 + \frac12e_0$,
$e_1 = 1 + \frac12e_2 + \frac12e_0$ (after $H$, a $T$ destroys progress), $e_2 = 1 + \frac12e_2 + 0$ (after $HH$, a further $H$ keeps $HH$); so $e_2 = 2$, $e_0 = 2 + e_1$, $e_1 = 2 + \frac12e_0$, and $e_0 = 8$. The difference comes from
the self-overlap of $HTH$.
\end{solution}

\begin{exercise}\label{martingales:exercise-polya}
In P\'olya's urn (Example~\ref{martingales:example-convergence}(1)), show that the number of red balls after $n$ draws is uniformly distributed on $\{1, \ldots, n + 1\}$, and deduce that the limit of $R_n$ is uniform on $[0, 1]$.
\end{exercise}

\begin{solution}
A given sequence of draws with $k$ reds and $n - k$ blues has probability $\frac{k!\,(n - k)!}{(n + 1)!}$: the numerators of the successive probabilities are $1, 2, \ldots, k$ for the red draws and $1, \ldots, n - k$ for the blue ones,
in some order, and the denominators are $2, 3, \ldots, n + 1$. Multiplying by $\binom nk$ gives $\frac{1}{n + 1}$ for each $k \in \{0, \ldots, n\}$, that is $k + 1$ red balls. So $R_n$ is uniform on $\{\frac{1}{n + 2}, \ldots, \frac{n + 1}{n + 2}\}$, which
converges in distribution to the uniform law; since $R_n$ converges almost surely, the limit is uniform (Proposition~\ref{probability:proposition-convergence-distribution}).
\end{solution}

\begin{exercise}\label{martingales:exercise-walk-expected-time}
Let $S_n$ be the simple random walk started at $0$ and $T = \inf\{n : S_n = 1\}$. Show without computation that $E[T] = \infty$, using Wald's identity.
\end{exercise}

\begin{solution}
$T < \infty$ almost surely (Counterexample~\ref{martingales:counterexample-optional-stopping}). If $E[T] < \infty$, Theorem~\ref{martingales:theorem-wald} would give $E[S_T] = 0\cdot E[T] = 0$, but $S_T = 1$.
\end{solution}

\begin{exercise}\label{martingales:exercise-knight}
Show that the knight's graph on an $8 \times 8$ chessboard has $168$ edges, and deduce Example~\ref{martingales:example-knight}.
\end{exercise}

\begin{solution}
A knight move is a vector $\pm(1, 2)$, $\pm(2, 1)$, $\pm(1, -2)$ or $\pm(2, -1)$; up to sign there are $4$ directions. The move by $(1, 2)$ is possible from the $7 \cdot 6 = 42$ squares $(x, y)$ with $x \leq 7$, $y \leq 6$, and likewise $42$
for each other direction, so there are $4 \cdot 42 = 168$ edges and the total degree is $336$. A corner has degree $2$; by Theorem~\ref{martingales:theorem-stationary} and Example~\ref{martingales:example-stationary}(1),
$E_{\text{corner}}[T] = 336/2 = 168$ (the graph is connected).
\end{solution}

\begin{exercise}\label{martingales:exercise-branching-poisson}
For a branching process with $\mathrm{Poi}(m)$ offspring, show that the extinction probability $q$ satisfies $q = e^{m(q - 1)}$, and compute $q$ numerically for $m = 2$.
\end{exercise}

\begin{solution}
$f(s) = \sum_ke^{-m}\frac{m^ks^k}{k!} = e^{m(s - 1)}$, and $q$ is the smallest fixed point by Theorem~\ref{martingales:theorem-branching}. For $m = 2$, iterating $q_{n+1} = e^{2(q_n - 1)}$ from $q_0 = 0$ gives
$0.135, 0.177, 0.193, 0.199, \ldots \to q \approx 0.2032$.
\end{solution}

\begin{exercise}\label{martingales:exercise-wiener-integral}
For a deterministic $f \in L^2([0, \infty))$, show that $\int_0^\infty f(s)\,dB_s$ has law $N(0, \|f\|_2^2)$. Deduce that $\int_0^ts\,dB_s \sim N(0, \frac{t^3}{3})$.
\end{exercise}

\begin{solution}
For a step function $f = \sum_kc_k\mathbf{1}_{(t_k, t_{k+1}]}$, the integral $\sum_kc_k\Delta_kB$ is a sum of independent normal variables, hence $N(0, \sum c_k^2(t_{k+1} - t_k)) = N(0, \|f\|_2^2)$
(Example~\ref{probability:example-characteristic-functions}(2)). Step functions are dense in $L^2$, the integral is an isometry (Theorem~\ref{martingales:theorem-ito-integral}), and $L^2$ limits of $N(0, \sigma_n^2)$ variables with
$\sigma_n \to \sigma$ are $N(0, \sigma^2)$: their characteristic functions converge (by $|e^{ix} - e^{iy}| \leq |x - y|$) to $e^{-\sigma^2t^2/2}$. For $f(s) = s\mathbf{1}_{[0, t]}$, $\|f\|_2^2 = t^3/3$.
\end{solution}

\begin{exercise}\label{martingales:exercise-ehrenfest-return}
For the Ehrenfest urn with $N$ molecules, show that the expected return time to the state where all molecules are on the left is $2^N$, while the expected return time to $N/2$ ($N$ even) is $2^N/\binom{N}{N/2} \sim \sqrt{\pi N/2}$.
\end{exercise}

\begin{solution}
The chain is irreducible on a finite set with reversible distribution $\pi(k) = \binom Nk2^{-N}$ (Example~\ref{martingales:example-stationary}(2)); by Theorem~\ref{martingales:theorem-stationary}, $E_k[T_k] = 1/\pi(k)$. For $k = N$ this is
$2^N$; for $k = N/2$ it is $2^N/\binom N{N/2}$, and $\binom{N}{N/2} \sim 2^N\sqrt{2/(\pi N)}$ by Stirling's formula. This resolves the paradox of Loschmidt and Zermelo: the second law of thermodynamics is not violated in
practice, since for $N = 10^{23}$ the recurrence time $2^N$ dwarfs the age of the universe.
\end{solution}
